% sec:kerr-bl-quartic — The Boyer-Lindquist Radius Quartic
\section{The Boyer-Lindquist Radius Quartic}
\label{sec:kerr-bl-quartic}

The Kerr--Schild metric~\cref{eq:kerr-ks-metric} is expressed in
Cartesian coordinates $(\bar{t}, x, y, z)$, but the scalar function
$f = 2Mr/\Sigma$ and the null covector $l_\mu$ both depend on the
Boyer--Lindquist radius $r$, which is not a coordinate in the
Kerr--Schild chart.  Every evaluation of the metric --- and therefore
every step of the geodesic integrator --- requires extracting $r$
from the Cartesian position.  The relationship between $r$ and
$(x, y, z)$ is algebraic but nonlinear, governed by a quartic whose
derivation and closed-form solution are the subject of this section.

The key geometric idea is that surfaces of constant Boyer--Lindquist
radius are not spheres but oblate ellipsoids, flattened along the
spin axis by the parameter $a$.  A Cartesian point $(x, y, z)$ lies
on exactly one such ellipsoid, and identifying which one amounts to
solving a polynomial in $r$.  The polynomial turns out to be quartic,
but its special structure --- only even powers of $r$ appear ---
reduces it to a quadratic in $r^2$ with a closed-form solution.
\coderef{Spacetime.Kerr}{blRadius}

\paragraph{Derivation.}

The oblate spheroidal relations~\cref{eq:kerr-oblate-spheroidal}
connect the Kerr--Schild Cartesian coordinates to the
Boyer--Lindquist radius $r$ and polar angle $\theta$.  Taking the
squared modulus of $x + iy = (r + ia)\,e^{i\bar{\varphi}}\sin\theta$
gives
%
\begin{equation}
  x^2 + y^2 = (r^2 + a^2)\sin^2\!\theta \,,
  \label{eq:bl-xy-relation}
\end{equation}
%
and the second relation $z = r\cos\theta$ yields
$\cos^2\!\theta = z^2/r^2$, hence
$\sin^2\!\theta = 1 - z^2/r^2$.  Writing $p^2 \equiv x^2 + y^2 + z^2$
for the squared Euclidean distance from the origin (unrelated to the
pressure $p$ of \cref{sec:gr-stress-energy} or the four-momentum
$p_\mu$ of \cref{sec:sr-causal}; the codebase variable
\texttt{p2} uses the same symbol), we find
%
\begin{equation}
  p^2 = (r^2 + a^2)\!\left(1 - \frac{z^2}{r^2}\right) + z^2
      = r^2 + a^2 - \frac{a^2 z^2}{r^2} \,.
  \label{eq:bl-p2-relation}
\end{equation}
%
The first equality substitutes \cref{eq:bl-xy-relation} and $z^2$;
the second expands and collects terms.  Multiplying through by $r^2$
clears the denominator and yields the quartic stated in
\cref{sec:kerr-bl}:
%
\begin{equation}
  r^4 - (p^2 - a^2)\,r^2 - a^2 z^2 = 0 \,.
  \label{eq:bl-quartic-derived}
\end{equation}
%
This is the implicit definition of the Boyer--Lindquist radius as a
function of Cartesian position.  Despite being quartic in $r$, the
equation contains only even powers --- it is biquadratic, and
therefore reducible to a quadratic in $r^2$.

\paragraph{Closed-form solution.}

Writing $u = r^2$, the quartic becomes
%
\begin{equation}
  u^2 - (p^2 - a^2)\,u - a^2 z^2 = 0 \,.
  \label{eq:bl-quadratic-u}
\end{equation}
%
The quadratic formula gives two roots:
%
\begin{equation}
  u_\pm = \frac{(p^2 - a^2) \pm
    \sqrt{(p^2 - a^2)^2 + 4a^2 z^2}}{2} \,.
  \label{eq:bl-u-roots}
\end{equation}
%
The discriminant $\mathcal{D} = (p^2 - a^2)^2 + 4a^2 z^2$ is a sum
of two squares --- $(p^2 - a^2)^2$ and $(2az)^2$ --- and therefore
non-negative for all $(x, y, z)$, guaranteeing real roots.  Since
$\sqrt{\mathcal{D}} \geq \abs{p^2 - a^2}$, the smaller root
satisfies
$u_- = \tfrac{1}{2}\bigl[(p^2 - a^2) - \sqrt{\mathcal{D}}\bigr]
\leq \tfrac{1}{2}\bigl[(p^2 - a^2) - \abs{p^2 - a^2}\bigr] \leq 0$,
with equality only at the ring singularity ($p = \abs{a}$, $z = 0$)
where both roots vanish.  Only the larger root $u_+$ yields a
positive value of $r^2$, so the Boyer--Lindquist radius is uniquely
determined:
%
\begin{equation}
  r = \sqrt{u_+}
    = \sqrt{\frac{(p^2 - a^2)
      + \sqrt{(p^2 - a^2)^2 + 4a^2 z^2}}{2}} \,.
  \label{eq:bl-r-solution}
\end{equation}
%
The expression involves one square root nested inside another, but no
transcendental functions --- the cost is comparable to a handful of
multiplications, making it negligible relative to the full metric
evaluation.
\implnote{The codebase variable \texttt{disc} in \texttt{blRadius}
stores $\mathcal{D}$; \texttt{diff} stores $p^2 - a^2$.  The nested
square roots appear as \texttt{sqrt r2} where
\texttt{r2 = (diff + sqrt disc) / 2}.}

\paragraph{Limiting cases.}

Setting $a = 0$ collapses the discriminant to
$\mathcal{D} = p^4$ and the solution to $r = \sqrt{p^2} = p$: the
Boyer--Lindquist radius equals the Euclidean distance, confirming
the Schwarzschild limit.  The codebase guards this case with an
exact-equality branch in \texttt{blRadius}
(\cref{sec:kerr-bl}), avoiding the unnecessary discriminant
computation.

In the equatorial plane ($z = 0$), the discriminant reduces to
$(p^2 - a^2)^2$ and therefore
$\sqrt{\mathcal{D}} = \abs{p^2 - a^2}$, giving
$u_+ = \tfrac{1}{2}\bigl[(p^2-a^2) + \abs{p^2-a^2}\bigr]
= \max(p^2 - a^2,\; 0)$.
The solution simplifies to
%
\begin{equation}
  r\big|_{z=0} = \sqrt{\max(p^2 - a^2,\; 0)} \,.
  \label{eq:bl-equatorial}
\end{equation}
%
The BL radius vanishes when $p^2 = a^2$ and $z = 0$,
i.e.\ when $x^2 + y^2 = a^2$ --- a circle of Euclidean radius
$\abs{a}$ in the equatorial plane.  This is the ring singularity at
$r = 0$, $\theta = \pi/2$, where $\Sigma = 0$ and the curvature
diverges (\cref{sec:kerr-bl}).
\physnote{The ring singularity lies at Cartesian distance $\abs{a}$
from the spin axis in the Kerr--Schild embedding.  For a
$10\,M_\odot$ black hole with $a/M = 0.9$, this Euclidean radius is
$\abs{a} \approx 13\;\text{km}$.}

Along the spin axis ($x = y = 0$), we have $p^2 = z^2$, so
$\mathcal{D} = (z^2 - a^2)^2 + 4a^2 z^2 = (z^2 + a^2)^2$, giving
$r = \sqrt{z^2} = \abs{z}$.  On the axis, the BL and Euclidean
radii coincide regardless of the spin --- the oblate spheroidal
distortion vanishes at the poles.

These limits serve as consistency checks: the formula recovers
ordinary spherical geometry when spin vanishes, and localises the
coordinate pathology exactly where the curvature is known to diverge.

\paragraph{Numerical robustness.}

The closed-form solution \cref{eq:bl-r-solution} is numerically
well-behaved everywhere except in the immediate neighbourhood of
the ring singularity, where $r \to 0$ while $f \to \infty$.
Because $\mathcal{D}$ is a sum of non-negative terms, the
floating-point subtraction that causes catastrophic cancellation in
expressions of the form $A - B$ cannot arise in its computation.
The only subtlety is the outer square root: when $p^2 \gg a^2$, the
expression $u_+ \approx p^2 - a^2 z^2/p^2$ is dominated by $p^2$,
and expanding both square roots yields $r \approx p$ with relative
error $\order{a^2/p^2}$.  No iterative refinement or
special-function evaluation is needed; the codebase computes $r$
with a fixed number of floating-point operations at every
integration step.
\warnnote{The quartic \cref{eq:bl-quartic-derived} has no solution
with $r^2 > 0$ when $\mathcal{D} < 0$, but this cannot occur since
$\mathcal{D}$ is a sum of squares.  Any negative discriminant in a
numerical implementation indicates a programming error, not a
physical edge case.}
